\documentclass[10pt]{article}
\usepackage[margin=20mm]{geometry}
\usepackage{amsmath,amssymb,mathtools,fontspec,unicode-math,xurl,hyperref}
\IfFontExistsTF{DejaVu Serif}{\setmainfont{DejaVu Serif}}{\setmainfont{Latin Modern Roman}}
\setmathfont{latinmodern-math.otf}
\setlength{\parindent}{0pt}\setlength{\parskip}{6pt}\sloppy
\title{Technical review of
On the ergodicity of infinite antisymmetric extensions of symmetric IETs}
\author{AI technical verification}
\date{}
\begin{document}
\maketitle
Abstract: This review concerns the 19-page arXiv version 2304.01868v4 by Przemysław Berk and Frank Trujillo. All source pages were read, with focused checks of essential values, reflection identities, Rauzy towers, the centre-projection input and the final rigidity construction. A central-tower branch doubles the same return tower, invalidating the local Birkhoff-sum calculation used to produce the essential value minus one. A positive twelve-step Rauzy example, checked by exact arithmetic, makes the error explicit. This is a Major defect in the presented proof; it does not establish that the main ergodicity theorem is false. Two local statement corrections concern positive return times and the order of quantifiers. Infinite subsequence selection and a complete repair of the main proof remain outside the certified result. Scientific adjudication remains pending.\par
\section*{Source identity and scope}
The assigned identifier is M25-0262. The frozen title is On the ergodicity of infinite antisymmetric extensions of symmetric IETs; its authors are Przemysław Berk and Frank Trujillo, and the journal DOI is `10.4171/jems/1635`. The actual artifact is arXiv `2304.01868v4`, dated 4 April 2024 (submission time 14:39:39 UTC), with 19 pages. The 2025 batch membership is retained. No equivalence to a later publisher version is assumed. All page references below are physical PDF pages, which agree with the printed pagination in this manuscript.\par
The complete source, including its bibliography, was read. The independent checks concern the hypotheses of Theorem 2.1, the definitions and criterion in Section 4, the cancellation identities in Section 5, the quantifiers and centre locations in Lemma 6.1, and the length, height, measure and Birkhoff-sum steps of Section 7. Selected source pages were inspected as images. The consulted part of the cited centre-location paper [BTU] is pages 19–21, including Lemma 3.12 and its proof; this is not a claim to have independently reviewed that entire paper.\par
The manifest’s bounded searches identified no supplementary mathematical file or author code. No author code was executed. A separate correction search did not establish an applicable correction to this exact version; the inconclusive search is not evidence that none exists.\par
\section*{Framework and checks that survive}
For a symmetric interval exchange \(T\) on \(I=[0,1)\), the cocycle is \[f=\chi_{(0,1/2)}-\chi_{(1/2,1)},\qquad
 T_f(x,m)=(Tx,m+f(x)).\] Theorem 2.1 asserts ergodicity of \(T_f\) for Lebesgue almost every length parameter, with product Lebesgue and counting measure. Its quantifier does not assert the conclusion for every minimal IET. The special coboundary example and the parity obstruction to ergodicity of \(T_f\times T_f\) therefore do not contradict the theorem.\par
Away from the exceptional discontinuity orbits, reflection gives \(\iota T=T^{-1}\iota\), and the negative cocycle convention is \(S_{-n}f=-S_nf\circ T^{-n}\). These give the identities in Lemma 5.2 and Corollary 5.3. In particular, \[S_{2n}f(T^{-n}(1/2))=-f(T^n(1/2)).\] This sign was independently checked; it is not the cause of Problem 1.\par
The partial-rigidity criterion of Proposition 4.2 has a direct verification. If \(\mu(\Xi_j)\to\delta>0\) and \(\mu(T\Xi_j\mathbin{\triangle}\Xi_j)\to0\), every weak-star limit of \(\chi_{\Xi_j}\) is invariant and hence equals \(\delta\) by ergodicity. Approximation of \(\chi_A\) by continuous functions on the compact metric space, together with uniform closeness of \(T^{h_j}x\) to \(x\) on \(\Xi_j\), then gives positive measure to \(A\cap T^{-h_j}A\cap\Xi_j\). If \(S_{h_j}f=a\) there, this supplies the required positive-time essential-value witness. Thus the outstanding core difficulty is the claimed fixed value on the constructed sets.\par
The measure estimate in the odd case uses \[\mathop{\mathrm{Leb}}(\Xi_k)=q_k\lambda_\alpha^{n_k}/5.\] The preceding sentence on page 14 omits the factor \(1/5\), but the subsequent constant \(10d\rho(\gamma)\) uses the correct factor. Likewise the length comparisons must be read homogeneously after normalizing by \(|I^{n_k}|\). These local presentation slips do not by themselves invalidate the positive lower bound. The near-invariance estimate is \(\mu(T\Xi_k\triangle\Xi_k)\leq2\mu(B_k)\to0\). Neither fact establishes the Birkhoff-sum identity.\par
\section*{Confirmed findings}
Problem 1 (The central branch counts the same tower twice).  In Case A, the time \(h_k=h_\alpha^{n_k}+h_{\bar\alpha}^{n_k}\) is defined on page 12. Figure 1(ii), page 15, explicitly includes \(\alpha=\bar\alpha\), the central label. Claim A4, page 14, asserts \[S_{h_k}f(x)-S_{h_k}f(1/2)=f(T^ix),\qquad 0\leq i<q_k. \tag{A4}\] Its proof counts each level of the two towers once. When the two labels coincide, these are the same tower, traversed twice. The asserted calculation is false in this permitted local branch.\par
\subsection*{A genuine positive Rauzy configuration}
Use labels \(a,b,c\) and the symmetric permutation \[\pi=\begin{pmatrix}a&b&c\\c&b&a\end{pmatrix}.\] Let \(t\) and \(b\) denote top and bottom wins, respectively. The word \(\eta=(tbtbtb)^2\) returns to \(\pi\) and has length matrix \[B=\begin{pmatrix}5&8&4\\12&21&8\\12&20&9\end{pmatrix},
 \qquad \det B=1.\] All entries are positive. Define \[\mu_a=1+\frac{\sqrt2}{10^5},\quad
 \mu_b=1,\quad
 \mu_c=1-\frac{\sqrt3}{10^5},\qquad
 L=\mathbf1^{\top}B\mu,\quad \lambda=B\mu/L.\] Every winner inequality in the twelve-step path is strict. The induced lengths at time 12 are \(\mu/L\), the induced permutation is \(\pi\), and the tower heights are the column sums \[(h_a^{12},h_b^{12},h_c^{12})=(29,49,21).\] The normalized vector \(\mu/|\mu|\) belongs to the open set in equation (6) on page 10: each coordinate differs from \(1/3\) by less than \(1/2700\).\par
The three coordinates of \(\mu\) are linearly independent over \(\mathbb Q\). The invertible integer matrix \(B\) preserves this property. Thus this is not a zero-length or terminating rational example. A connection between an interior discontinuity and another would impose a nonzero integer relation on the lengths. The same affine calculation, with half-integer centre coefficients, excludes the midpoint’s landing on a discontinuity in the checked orbit.\par
Let \(p_{12}\) be the first backwards projection to \(I^{12}\). Direct iteration gives \[\begin{array}{c|rrrr}
 z&c_a&c_b&c_c&1/2\\ \hline
 \text{projection time}&17&27&13&3\\
 p_{12}(z)&c_a^{12}&c_b^{12}&c_c^{12}&|I^{12}|/2 .
\end{array}\] Consequently the centre-set identity used from Lemma 6.1 holds at this time, \(\ell=3\), \(\alpha=\bar\alpha=b\), and \(q=21\). The paper’s time is \(h=98\).\par
\subsection*{Exact violation}
Equation reference: eq:sums. Set \[y=c_b^{12}+\frac{3\mu_b}{10L},\qquad x=T^3y.\] Since \(1/5<3/10<2/5\), this is a point of the specified base \(B_k\); take \(i=0\) in (A4). In particular, \[x-\frac12=\frac{(\mu_a-\mu_c)/2+3\mu_b/10}{L}>0,\qquad f(x)=-1.\] Exact iteration of the original IET gives \[\begin{aligned}
 S_{49}f(1/2)&=0,& S_{98}f(1/2)&=1,\\
 S_{49}f(x)&=-1,& S_{98}f(x)&=-2.
 \end{aligned}\] Thus the two sides of (A4) are \(-3\) and \(-1\). Equation (15), page 15, asserting \(S_hf(1/2)=0\), also fails here.\par
For transparency, the midpoint’s first 49 terms contain 24 positive values, 24 negative values and its initial zero. The next 49 contain 25 positive and 24 negative values. The 98 terms at \(x\) contain 48 positive and 50 negative values. The independent calculation stores each Rauzy step, each centre projection, the sign lists, and 2,039 certified comparisons. Coordinates are represented as \(a+b\sqrt2+c\sqrt3\) with rational coefficients. Radical signs are certified by rational lower and upper bounds obtained with integer square roots; no floating-point branch decision or author code is used.\par
There is also a short obstruction independent of these particular sums. In the central branch \(h=2h_\alpha\) is even. Off the null set of orbits meeting \(0\) or \(1/2\), every term of \(S_hf\) is \(\pm1\), and hence \(S_hf\in2\mathbb Z\). Equation (13), which requires the odd value \(-1\) on \(\Xi_k\), cannot hold almost everywhere in this branch.\par
\subsection*{Scope, counterevidence and impact}
The finite example checks the positive path, symmetric return, balanced induced lengths, centre projections, tower heights and the admissible point used by the local argument. Its branch calculations persist for nearby lengths; they are not caused by a boundary convention. A bounded affine check gives strict partition margins at \(\mu=(1,1,1)\) of at least \(0.5\) for the midpoint and \(0.2\) for \(x\). The relevant coefficient norms are at most \(40.5\) and \(40.8\). Consequently the inequalities persist in \(|\mu_j-1|<10^{-4}\), \(\mu_a>\mu_c\); the sign at time 49 is set by \(T^{49}(1/2)-1/2=(\mu_c-\mu_a)/L<0\).\par
This review has not constructed a whole infinite sequence with the specified fixed suffix and time 12 among its members. Nor has it proved that every selectable good subsequence has infinitely many central times. The counterexample refutes the local calculation from the local hypotheses actually used, and the parity argument applies whenever the explicitly admitted central branch occurs. Lemma 6.1 and the consulted centre-location lemma do not state an exclusion of that branch. Avoiding it would require an additional selection argument.\par
The dependency is Claim A4 \(\Rightarrow\) equations (14)–(15) \(\Rightarrow\) equation (13) \(\Rightarrow\) the use of Proposition 4.2 to obtain \(-1\in\mathop{\mathrm{Ess}}(T_f)\) in the odd case. This is a central derivation, so the assessment is Major. The assessment neither requires nor asserts a counterexample to Theorem 2.1 itself.\par
\subsection*{Repair boundary}
A concrete repair direction is to use \(h_\alpha\) once when \(\alpha=\bar\alpha\), retaining \(h_\alpha+h_{\bar\alpha}\) otherwise. The displacement estimate (9) is then stronger than the needed central-interval margin, and a single passage restores the once-per-level count. The two values at time 49 in [eq:sums] (\#eq:sums) verify this local repair on the explicit example. Combining the corrected single-passage identity with Corollary 5.3 gives the intended zero at the midpoint under the stated local geometry. A complete revision must still check the global choice of times, the partial-rigidity estimates and all branch dependencies; this report does not certify that revision as completed or adopted by the authors.\par
Problem 2 (Positive-time essential values do not satisfy the stated general theorem).  The definition on page 7 requires \(n\in\mathbb N\), but Theorem 4.1 then asserts a subgroup property and the full ergodicity criterion for arbitrary measurable cocycles.\par
Take a one-point probability space, \(T=\mathrm{id}\), \(G=\mathbb Z\) with counting Haar measure and \(\varphi=1\). This meets every displayed hypothesis of Theorem 4.1. For a sufficiently small ball around an integer \(a\), the definition holds precisely when \(a=n\) for an allowed \(n\in\mathbb N\). Hence \(\mathop{\mathrm{Ess}}(T_\varphi)\) is the positive integers (or the nonnegative integers if \(0\in\mathbb N\)), not a subgroup. Meanwhile \(T_\varphi\) is the translation \(m\mapsto m+1\) on \(\mathbb Z\); every invariant set is empty or all of \(\mathbb Z\), so it is ergodic. It therefore also contradicts the asserted equivalence.\par
A consistent invertible formulation uses integer times, including the identity time in the usual convention, or states appropriate recurrence hypotheses and proves the positive-time version. The main cocycle is zero-mean, and Proposition 4.2 supplies the specific positive-time witness needed there, so this example alone establishes no additional failure of the main conclusion. The assessment is Minor, a local definition–theorem mismatch. The recurrence assumption must not be silently added to the theorem when assessing its present text.\par
Problem 3 (The literal order of quantifiers in Lemma 6.1 is too strong).  Page 9 says that for almost every symmetric IET \(T\), and for every positive-measure subset \(A\) of the length simplex, certain normalized Rauzy iterates belong to \(A\).\par
Fix any \(T\) for which induction is indefinitely defined, and write \[D_T=\{\lambda^n/|\lambda^n|:n\geq0\},\qquad A_T=\Lambda_{\mathcal A}
 \setminus D_T .\] For \(d\geq2\), \(D_T\) is countable and Borel, while \(A_T\) has full Lebesgue measure in the simplex. No normalized iterate belongs to \(A_T\). Thus no single \(T\) can satisfy the literal universal assertion over all positive-measure measurable sets.\par
The needed correction is: for each fixed measurable \(A\) of positive measure, the assertion holds for almost every \(T\), allowing the null exceptional set to depend on \(A\). Alternatively, a separately proved uniform statement for open sets may suffice. Section 7 fixes its open set \(A\) in advance for each \(d\), so this correction preserves the intended use; it does not itself repair Problem 1. The consulted recurrence proof in [BTU], pages 19–21, chooses a fixed open neighbourhood before applying recurrence, consistent with this distinction. Assessment: Minor, at the level of the literal statement.\par
\section*{Coverage limits and disposition}
The even case was read, including its special three-tower time and Claims B1–B4. Its measure and invariance estimates were checked, but its entire multi-visit cancellation argument was not independently reconstructed. A possible analogous counting difficulty was not promoted to a finding without a complete admissible counterexample. Boundary identities outside the generic discontinuity-free orbit were likewise not used to claim a main-theorem defect.\par
The foundational Rauzy–Veech ergodicity results, all external cocycle criteria and a full infinite central-return construction were not independently reproved. These are explicit limits on the coverage, not a substitute for a counterexample. The assessment is one confirmed Major proof-step issue and two local Minor statement issues. The current report concerns only the exact v4 source. Author adoption is unknown; no human review or independent acceptance is asserted.\par
9 P. Berk and F. Trujillo, On the ergodicity of infinite antisymmetric extensions of symmetric IETs, arXiv:2304.01868v4, 4 April 2024, 19 pp. Frozen journal DOI: 10.4171/jems/1635. P. Berk, F. Trujillo and C. Ulcigrai, Ergodicity of explicit logarithmic cocycles over IETs, arXiv:2210.16343v2, 4 August 2023. Consulted pages 19–21 only.\par
\end{document}
